Data memory implementation depends on the chip variant:

\begin{itemize}
  \item \textbf{Variant~A (No-Cache):}
    An 8\,kB on-chip SRAM (\texttt{asDMemTop}) serves as the data memory.
    It is connected to the CPU via a Wishbone B3 slave interface and
    supports byte, halfword, word, and doubleword accesses.
    Address range: \texttt{0x0000\_0000}--\texttt{0x0000\_1FFF}.

  \item \textbf{Variant~B (Cache):}
    There is no traditional on-chip data memory.
    The data address space is split into two regions:
    \begin{itemize}
      \item \textbf{Flash region} (addresses below the Scratchpad base):
        read-only data cached in the 4\,kB D-Cache.
        Only load misses trigger cache-line fills from NOR Flash.
        Store operations to this region are not permitted and generate a
        bus error (\texttt{dc\_err\_o}).
      \item \textbf{Scratchpad region}
        (\texttt{SP\_BASE}~=~\texttt{0x0000\_0000},
        size~=~8\,kB by default):
        an on-chip SRAM accessed directly with byte granularity,
        bypassing the D-Cache.
        This region holds \texttt{.data}, \texttt{.bss}, heap, and stack.
    \end{itemize}
\end{itemize}
